% Section 3: Data and setup.
% All numbers come from macros in 03_data_numbers.tex and the table in
% ../tables/tab_data_counts.tex, both generated by scripts/05_section3_numbers.py
% from committed results/. Do not type numbers into this file.
% Citations: sisbar_refs.bib (You_2023, Wolock_2019).
% AUTO-GENERATED by scripts/05_section3_numbers.py -- do not edit by hand.
\newcommand{\nLibraries}{22}
\newcommand{\nGenes}{33,694}
\newcommand{\nBarcodes}{215,960}
\newcommand{\nPublishedCells}{173,564}
\newcommand{\nCloneTableCells}{115,046}
\newcommand{\pctMultiBarcode}{49}
\newcommand{\pctMultiBarcodeMin}{36}
\newcommand{\pctMultiBarcodeMax}{53}
\newcommand{\nStageIVCells}{5,329}
\newcommand{\nStageIVCollisions}{20}
\newcommand{\nExpThreeClones}{555}
\newcommand{\nExpThreeCells}{2,692}
\newcommand{\nExpThreePubClones}{559}
\newcommand{\nExpThreePubCells}{2,816}
\newcommand{\pctExpThreeCloneGap}{0.7}
\newcommand{\pctExpThreeCellGap}{4.4}
\newcommand{\swapRepTwoGSM}{GSM6889993, GSM6889994}
\newcommand{\swapRepThreeGSM}{GSM6889987, GSM6889988}
\newcommand{\nPaperQCPass}{178,226}
\newcommand{\nCloneTableFailPaperQC}{11,350}
\newcommand{\pctCloneTableFailPaperQC}{9.9}
\newcommand{\nCloneTableOverUMICap}{10,973}
\newcommand{\pctStageIVOverUMICap}{40}
\newcommand{\nExpThreeClonesPaperQC}{437}
\newcommand{\umiCap}{30,000}
\newcommand{\nPrimaryCells}{198,073}
\newcommand{\nCloneTablePrimary}{114,667}
\newcommand{\nCloneTableDropped}{379}
\newcommand{\pctCloneTableKept}{99.7}
\newcommand{\nScrubletAutoTotal}{5,896}
\newcommand{\nScrubletAutoMax}{2,434}
\newcommand{\nLibsScrubletZero}{12}
\newcommand{\rhoDoubletUMI}{0.18}
\newcommand{\umiRatioFlagMin}{1.7}
\newcommand{\umiRatioFlagMax}{2.9}
\newcommand{\nSensitivityCells}{181,018}
\newcommand{\nCloneTableSensitivity}{103,049}
\newcommand{\nDACellsCrossStage}{8}
\newcommand{\minGenes}{1,000}
\newcommand{\maxMt}{15}
\newcommand{\dblRatePct}{0.8}
\newcommand{\dblRatePer}{1,000}
\newcommand{\dayOne}{13}
\newcommand{\dayTwo}{21}
\newcommand{\dayThree}{30}
\newcommand{\dayFour}{45}


\section{Data and Setup}\label{sec:data}

\subsection{Dataset and lineage design}
We use the single-cell split-barcoding (SISBAR) dataset of You \emph{et al.}~\cite{You_2023},
in which H9 human embryonic stem cells are differentiated towards ventral midbrain and
hindbrain fates and profiled by droplet single-cell RNA sequencing at four stages:
I (day~\dayOne), II (day~\dayTwo), III (day~\dayThree) and IV (day~\dayFour). In each of
three experiments, progenitors are labelled with a retroviral barcode library at one
stage. After three to four divisions the culture is split: one half is profiled at that
stage and the other half at the next stage. Experiments~1, 2 and~3 therefore link
stages I$\to$II, II$\to$III and III$\to$IV. A barcode seen in both halves links
\emph{sister} cells, not the same cell observed twice, and this shared ancestry is the
ground truth used throughout. Each experiment is barcoded once and sampled at two
consecutive stages, so no clone spans three stages and every analysis is a single
transition $t\to t+1$.

Each experiment has three biological replicates. Replicate~1 is deposited as GEO
GSE207921 and replicates~2--3 as GSE221592, as \nLibraries{} 10x libraries holding
\nBarcodes{} cell barcodes and \nGenes{} genes. The authors' barcode-to-cell table
(GSE221592 supplementary file) lists \nCloneTableCells{} cells with their viral
barcodes and cluster labels.

\subsection{Clone definition and reproduction of published counts}
Of the clone-table cells, \pctMultiBarcode\% carry more than one viral barcode
(\pctMultiBarcodeMin--\pctMultiBarcodeMax\% across the six stage--experiment sheets),
reflecting multiple integrations per founder cell. We define a clone as the exact set of
viral barcodes carried by a cell; barcodes are specific to an experiment and replicate.
A clone is \emph{cross-stage} if it contains at least one cell at each stage of its
transition. Two alternative definitions, each barcode as its own clone and the
connected components of co-occurring barcodes, did not reproduce the published counts.
The barcode-set definition reproduces I$\to$II and II$\to$III exactly
(Table~\ref{tab:data_counts}). For III$\to$IV it yields \nExpThreeClones{} clones
containing \nExpThreeCells{} cells, against \nExpThreePubClones{} and
\nExpThreePubCells{} published (\pctExpThreeCloneGap\% fewer clones,
\pctExpThreeCellGap\% fewer cells). All \nStageIVCells{} stage-IV clone-table cells are
present in their expression matrix (Section~\ref{sec:data_libs}), so the shortfall lies in
the deposited table rather than in our matching. The handling rule was fixed in advance:
one matching pass, then accept the deposited count. No filter was adjusted to recover
the published value.

\subsection{Library and replicate assignment}\label{sec:data_libs}
Cell barcodes in the clone table carry a library suffix, except in stage~IV. For every
group of cells defined by suffix and viral-barcode replicate, each cell was found in
exactly the 10x library that the suffix and replicate predict. In stage~IV, the
replicate is read from the viral barcode, which ends in a replicate and experiment
code; this also resolves the \nStageIVCollisions{} stage-IV cell barcodes that occur in
more than one stage-IV library. In experiment~3, the replicate labels in the GEO sample
names are swapped relative to the viral barcodes, consistently in both stages: viral
replicate~2 corresponds to the samples named replicate~3 (\swapRepTwoGSM) and viral
replicate~3 to those named replicate~2 (\swapRepThreeGSM). Clone pairing is unaffected,
and we use the viral-barcode replicate as the replicate identifier throughout.

\subsection{Quality control}
The original study states thresholds of at least \minGenes{} detected genes, at most
\umiCap{} UMIs and at most \maxMt\% mitochondrial reads. Applied to all barcodes,
\nPaperQCPass{} cells pass, but \nCloneTableFailPaperQC{} (\pctCloneTableFailPaperQC\%)
of the authors' own clone-table cells fail. Nearly all of them (\nCloneTableOverUMICap)
exceed the UMI cap, including \pctStageIVOverUMICap\% of stage-IV clone-table cells,
consistent with the larger transcriptomes of maturing neurons. Under the stated
thresholds, III$\to$IV would fall to \nExpThreeClonesPaperQC{} cross-stage clones
(Table~\ref{tab:data_counts}). The deposited lineage table was therefore evidently not
filtered with the UMI cap. Optimal-transport couplings are fitted on barcoded and
unbarcoded cells together, so the same rule must apply to both. Our primary rule keeps
cells with at least \minGenes{} genes and at most \maxMt\% mitochondrial reads, without
a UMI cap. It retains \nPrimaryCells{} cells, including \nCloneTablePrimary{} of the
\nCloneTableCells{} clone-table cells (\pctCloneTableKept\%). Neither rule reproduces
the published \nPublishedCells{} high-quality cells, which reflect further filtering in
the original pipeline that we do not replicate.

\subsection{Doublets and the sensitivity cell set}
Doublet detection with Scrublet~\cite{Wolock_2019} was pre-specified, per library, with
an expected doublet rate of \dblRatePct\% per \dblRatePer{} barcodes. Its automatic
threshold proved unstable here. It flagged between none and \nScrubletAutoMax{} cells per
library (none in \nLibsScrubletZero{} of \nLibraries{} libraries; \nScrubletAutoTotal{}
in total), and its output depended on implementation settings. Observed scores also
overlapped those of simulated doublets, as expected when cells within a stage form a
continuum. We therefore do not filter on Scrublet. Instead we record each cell's doublet
score and flag the top expected-rate fraction in each library. Flagged cells have
\umiRatioFlagMin--\umiRatioFlagMax{} times the median UMI count of unflagged cells,
while the score itself correlates only weakly with UMI count (Spearman
$\rho=\rhoDoubletUMI$), so the flag is not a size cap in disguise. Every analysis for
RQ1--RQ3 is repeated with flagged cells excluded (\nSensitivityCells{} cells;
\nCloneTableSensitivity{} clone-table cells). This sensitivity analysis is
pre-specified and reported alongside the primary result.

\subsection{Analysis cell sets}
Table~\ref{tab:data_counts} lists cross-stage clones and their cells for each
transition and cell set. All analyses use one-step pairs from the primary set. Midbrain
dopaminergic neurons are rare among lineage-traced cells: only \nDACellsCrossStage{}
stage-IV dopaminergic cells fall in cross-stage III$\to$IV clones in the primary set.
Dopaminergic fate is therefore studied through progenitor fate bias at stage~III, not as
a forecasting target.

% AUTO-GENERATED by scripts/05_section3_numbers.py -- do not edit by hand.
\begin{table}[!t]
\centering
\caption{Cross-stage clones / cells in those clones, per transition and cell set.}
\label{tab:data_counts}
\footnotesize
\setlength{\tabcolsep}{3pt}
\begin{tabular}{lrrr}
\hline
Cell set & I$\to$II & II$\to$III & III$\to$IV \\
 & d13$\to$d21 & d21$\to$d30 & d30$\to$d45 \\
\hline
Published~\cite{You_2023} & 768 / 2,762 & 1,213 / 5,856 & 559 / 2,816 \\
All clone-table cells & 768 / 2,762 & 1,213 / 5,856 & 555 / 2,692 \\
Paper thresholds (not used) & 706 / 2,462 & 1,099 / 5,010 & 437 / 1,989 \\
Primary (this study) & 756 / 2,696 & 1,208 / 5,817 & 549 / 2,667 \\
Doublet sensitivity & 689 / 2,364 & 1,172 / 5,532 & 519 / 2,459 \\
\hline
\end{tabular}
\end{table}

